\chapter{文献综述与理论基础}

本章围绕研究问题，按"业务背景 → 服务模式变化 → 服务质量评价 → 测量方法 → 改进路径 → 研究空白"的概念链条展开文献综述。全章共六节，依次梳理普惠金融与小微贷款、金融科技服务模式、服务质量与客户满意度、SERVQUAL/IPA/AHP评价方法、DMAIC服务流程改进五条研究脉络，最后通过文献述评明确现有研究的不足与本文的定位。\par

\section{2.1 普惠金融与小微贷款研究}

小微企业的融资约束问题在学术界已形成较为系统的认识。其成因可从信息不对称、抵押不足、征信缺失与交易成本偏高四个维度加以理解。其一，小微企业普遍缺乏规范的财务报表和可验证的信用记录，金融机构难以准确评估其还款能力与还款意愿，由此形成严重的事前信息不对称，这种信息不对称还会引发逆向选择：风险较低的企业因无法自证信用而退出市场，留在信贷池中的企业平均风险上升，迫使金融机构进一步提高利率或收紧授信条件。其二，小微企业可用于抵押的固定资产有限，而信用担保体系在覆盖面和运作效率上的不足，使得"抵押缺口"成为制约信贷扩张的刚性约束。其三，小微企业征信信息的碎片化和非标准化特征，导致传统以央行征信报告为核心的授信决策模式面临严重的"信息采集成本-收益不对称"，这种交易成本的结构性偏高进一步压缩了金融机构服务小微的内生动力。上述三方面成因相互叠加，共同构成了小微企业融资难、融资贵的系统性根源。\par

上述信息不对称引发信贷配给的经典理论根源，可追溯至Stiglitz和Weiss（1981）的信贷配给理论\cite{ref17}。该理论证明，在信息不完美条件下，银行提高利率不仅无法出清超额信贷需求，反而会通过逆向选择和道德风险效应恶化贷款组合质量，因此银行宁愿维持低于市场出清水平的利率并实施信贷配给。该理论为理解小微融资困境提供了根本性的经济学解释：信息不对称是信贷配给的内生原因，破解小微融资困境不能仅依靠利率工具或担保工具，还需要从降低信息不对称、改善服务效率的方向寻找工程管理层面的解决方案。这一点构成了本文聚焦服务质量而非仅关注融资可得性的理论依据。\par

在实证层面，王贤彬等（2025）基于小微支行设立的准自然实验，发现普惠小微信贷供给的扩张显著促进了小微企业的市场进入，这一效应在金融服务相对薄弱的地区尤为明显\cite{ref2}。该研究从供给侧验证了信贷可及性对小微主体活力的关键作用，但其研究视角侧重于机构进入与业务扩张，并未深入讨论信贷服务过程中的客户体验维度。\par

在普惠信贷供给机制层面，政策性担保与风险分担工具被广泛讨论。宋卓霖、赵涵（2024）系统考察了政策性担保在小微融资中的作用边界，指出担保机制虽能部分缓解抵押不足问题，但其效果受制于担保机构的风险偏好和银行的协作意愿\cite{ref18}。这意味着单一的供给端政策难以独立破解小微融资困境，需要配套的信息治理和服务机制跟进。从方法论角度看，该研究提供了一个重要的启示：政策工具的实际效果高度依赖于执行层面的制度设计和流程配套，这恰是工程管理视角的切入点。\par

数字普惠金融的兴起为上述困局提供了新的解决思路。祝树金、朱捷尧（2025）的实证研究表明，数字普惠金融通过降低信息搜寻成本、拓展服务半径和加快资金流转，显著促进了中小微企业在价值链中的攀升\cite{ref7}。顾锋娟等（2024）利用省级面板数据进一步揭示，数字普惠金融对小微融资的促进效应存在地区异质性，东部地区的改善幅度大于中西部，这提示数字化红利的释放需要配套的基础设施和制度条件支撑\cite{ref21}。周雷等（2024）的研究则从数字普惠金融服务模式创新角度，探讨了如何通过场景化服务和信用信息共享缓解小微融资堵点\cite{ref22}。\par

国际研究同样提供了有力的对照证据。Yao和Yang（2022）基于中国创业板上市公司数据，发现数字金融通过缓解融资约束间接驱动了中小企业创新，其传导机制主要包括降低融资成本和缩短审批周期\cite{ref6}。Lu等（2021）的实证工作则聚焦于本地银行的作用，发现本地银行与数字金融的协同能够更有效地缓解中小企业的融资约束，这一"线上线下互补"的发现为理解服务渠道融合模式提供了实证基础\cite{ref13}。Zhang等（2023）的研究发现，数字普惠金融通过缓解融资约束促进中小企业技术创新，且这一效应在金融监管较强和政府补贴较充裕的环境下更为显著\cite{ref23}。\par

纵观上述文献，现有研究主要围绕"融资可得性"展开，即关注小微企业能否获得贷款、以什么成本获得贷款。这一视角虽然抓住了小微融资的核心矛盾，但在以下方面存在明显不足：贷款服务过程中的客户体验问题，包括申请便捷度、审批透明度、客户经理专业性和贷后服务质量等，尚未进入主流研究视野。对P公司这样的普惠金融机构而言，客户在申请、审批、放款、贷后各环节中感知到的服务效率、信息透明度和专业支持，同样构成了业务可持续性的关键变量。因此，本文主张将研究焦点从"融资可得性"适度扩展至"融资服务质量"，二者并非替代关系，而是互补关系。\par

\section{2.2 金融科技与小微贷款服务模式}

数字技术的快速渗透正在重塑小微贷款的服务模式。大数据风控使金融机构能够在缺乏传统抵押和财务报表的条件下，通过多维度替代数据（如经营流水、税务记录、线上行为轨迹）构建信用评估模型\cite{ref10}。汪勇等（2026）利用助贷机构数据证实，数字金融技术显著降低了借贷双方的信息不对称程度，但这种改善效果在客户数字素养较低或线上线下流程衔接不畅的场景中有所衰减\cite{ref10}。这一发现揭示了一个核心命题：技术赋能并非自动转化为服务体验提升，二者之间存在需要管理的"转化缺口"。\par

从服务模式演进的视角观察，小微贷款业务正经历从纯线下到线上线下融合、从纯人工到人机协同的转变。韩莉等（2021）在对金融科技助力小微企业融资的文献评析中，系统梳理了智能审批、线上获客、区块链征信等技术在小微信贷中的应用进展，同时指出一个关键缺失：多数研究关注的是技术本身的效率和精度，对技术嵌入后的客户体验变化着墨不多\cite{ref9}。Dumitrescu等（2021）将机器学习方法引入信用评分领域，证实了非线性决策树模型在预测精度上优于传统逻辑回归，但模型的可解释性下降可能带来新的客户信任问题\cite{ref24}。这一发现的管理含义在于：信用评分的技术进步如果以牺牲透明性为代价，可能在提升审批效率的同时损害客户对公平性的感知。\par

服务模式变革的另一面是"体验断点"的出现。线上流程虽然缩短了物理距离，但也可能拉大服务鸿沟。Davis（1989）提出的技术接受模型（TAM）以感知有用性和感知易用性为核心变量，系统解释了用户对信息技术系统的接受行为\cite{ref25}。Nugraha等（2022）针对印度尼西亚中小企业的调查发现，金融科技采纳意愿受到感知易用性和社会影响的显著驱动，但在实际使用中，用户对平台操作复杂、问题响应延迟、缺乏人工支持的抱怨较为普遍\cite{ref26}。相似的问题在中国语境下同样存在：线上申请虽然便捷，但当客户遇到资料补正、审批进度不明、退件原因笼统等情况时，缺乏有效的沟通渠道将加剧焦虑和不信任。吴勇民（2026）从金融大模型赋能供应链融资的视角，讨论了智能技术在融资服务中的双重效应：一方面提升了流程效率，另一方面也产生了新的协调成本和技术门槛\cite{ref27}。\par

对本研究而言，上述文献提供了两重启示。第一，数字技术是改善小微贷款服务质量的工具，不是目的。第二，技术引入与服务质量之间不存在天然的线性关系，需要系统性的服务设计和管理机制来弥合其中的落差。这恰好是DMAIC方法论的价值所在。\par

\section{2.3 服务质量与客户满意度研究}

服务质量（Service Quality）作为一个具有独立理论地位的管理概念，经历了从北欧学派到北美学派的概念化过程。Grönroos（1984）最早提出顾客感知服务质量模型，将服务质量区分为技术质量（what，即服务的结果）和功能质量（how，即服务的过程）。这一区分的价值在于：它指出了服务质量评价的双重维度，客户不仅关心"得到了什么"，也关心"是如何得到的"。Parasuraman、Zeithaml和Berry（1988）在此基础上发展出SERVQUAL量表，以期望-感知差距为核心逻辑，将服务质量操作化为可测量的多维度构念\cite{ref14}。Cronin和Taylor（1992）提出的SERVPERF模型则以绩效直接测量替代差距测量，认为仅测量感知绩效即可有效反映服务质量\cite{ref29}。该理论框架的基石在于一个根本性的认知：服务质量不是由服务提供者定义的，而是由客户通过比较事前期望与事后感知来判定的。这一认知范式对服务管理的颠覆性在于，它将质量定义权从供给侧转移到了需求侧。\par

在服务质量的结果变量方面，学术界积累了较为一致的经验证据。Nguyen等（2020）以银行业为背景，验证了"服务质量 → 客户满意度 → 转换成本 → 客户忠诚度"的完整关系链，其中服务质量对满意度的直接效应最为显著，而满意度在服务质量与忠诚度之间发挥部分中介作用\cite{ref31}。这一发现的管理含义是明确的：服务质量是可管理的前置变量，通过改进服务质量来提升满意度和忠诚度，比直接干预满意度或忠诚度更具操作可行性。在金融服务领域，该链条还获得了定量化的经验支持：服务质量每提升1分（5分制），客户忠诚度平均提升约0.6分，转换成本的提高进一步锁定了客户关系\cite{ref31}。Zaheer（2023）以银行业为对象，运用结构方程模型（SEM）进一步验证了SERVQUAL各维度对客户服务质量的差异化影响，为多维服务质量结构提供了统计证据\cite{ref32}。\par

上述关系链的理论基础可进一步追溯至Heskett等（1994）提出的服务利润链理论\cite{ref33}。该理论阐明了服务企业经营绩效的因果传导逻辑：内部服务质量驱动员工满意度，员工满意度驱动员工保留与生产率，进而创造外部服务价值；客户感知的服务价值驱动客户满意度，客户满意度驱动客户忠诚度，而客户忠诚度最终驱动收入增长和利润提升。在金融服务语境下，服务利润链将服务质量改进从成本中心重新界定为利润驱动要素：服务质量管理的目的不仅在于消除投诉，更在于构建可持续的客户关系和竞争优势，从而为普惠金融机构的服务质量投入提供了商业合理性论证。\par

银行业服务质量研究为本文提供了直接的参照系。Siddiqi（2011）针对孟加拉国零售银行的实证研究表明，服务质量各属性对客户满意度的影响存在显著差异，其中可靠性和响应性的贡献最大，有形性的贡献最小，这一发现与金融服务的无形性特征相符\cite{ref35}。Kheng等（2010）在马来西亚槟城银行的调查中得出了相似结论，并进一步发现服务质量对客户忠诚度的影响主要通过满意度的中介路径实现，这意味着即便服务质量得到改善，如果没有转化为客户可感知的满意度提升，其对忠诚度的拉动作用将大打折扣\cite{ref36}。在国内研究方面，孙彦佩（2025）以ZY银行郑州分行为案例，从网点服务流程的角度系统诊断了服务质量短板，其问题维度分类（人员服务、流程效率、信息沟通、环境设施）对本文第3章的现状诊断具有方法参考价值\cite{ref37}。黄兆锟（2025）聚焦于小微企业客户关系管理的优化，揭示了小微客户在贷款服务中对及时响应和个性化服务的高度敏感性，这种敏感性显著高于零售个人客户\cite{ref38}。\par

数字化转型为服务质量研究带来了新的变量。Zouari和Abdelhedi（2021）以伊斯兰银行为例，发现在数字化背景下，客户满意度的影响因素从传统的物理网点体验转向线上平台的可用性和安全性，这一转变的幅度在年轻客户群体中尤为显著\cite{ref39}。Kaur等（2021）针对印度北部数字银行的研究则提出了一个值得警惕的发现：数字化在提升便利性的同时，也引入了网络安全、数据隐私和服务中断等新型风险感知，这些风险感知对满意度存在显著的负向影响，其效应量在某些场景下甚至超过了便利性提升带来的正面收益\cite{ref40}。上述发现对本文的启示在于：小微贷款服务质量的评价维度不能简单照搬传统银行的量表，需要纳入便利性、透明性和风险安全等数字化时代的新维度。\par

综合本节梳理的文献，服务质量理论经过近四十年的积累，已经形成了"概念界定 → 维度操作化 → 前因后果验证"的成熟研究范式。在金融服务，尤其是小微贷款场景中，该范式的适用性已经得到初步验证，但针对小微贷款特定场景的维度改良和系统评价的工作仍然不足，这构成本文第4章评价模型的切入点。\par

\section{2.4 SERVQUAL、IPA与服务质量评价方法}

SERVQUAL模型是服务质量评价领域最具影响力的理论工具。Parasuraman、Zeithaml和Berry（1988）提出的经典五维模型将服务质量分解为有形性（Tangibles）、可靠性（Reliability）、响应性（Responsiveness）、保证性（Assurance）和移情性（Empathy），并通过期望-感知差距分值（Gap Score）来量化服务质量水平\cite{ref14}。该模型的核心优势在于：它将抽象的服务质量概念转化为可测量、可比较、可追踪的指标体系，为服务改进提供了量化起点。然而，SERVQUAL的通用性也是一把双刃剑，不同行业、不同服务场景需要根据自身特征对维度和指标进行酌情调整，这是模型应用中的基本共识。\par

在小微贷款服务场景中，经典五维的覆盖范围存在三个明显缺口。第一，数字渠道已成为小微贷款获客和受理的主要入口，便利性（如线上操作便捷度、材料一次告知率）直接影响客户的第一感知，但经典五维并未系统纳入。第二，小微贷款涉及复杂的费用结构和审批规则，透明性（如费用事先说明、进度实时可查）是建立客户信任的基础，同样是经典五维的盲区。第三，小微客户在经营周期、行业特征和资金需求上的异质性要求产品具备一定的适配弹性，这是服务针对性而非通用移情性的体现。因此，本文在第4章构建改良SERVQUAL模型时，将在经典五维基础上增加便利性、透明性和产品适配三个新维度，形成八维评价框架。\par

IPA（Importance-Performance Analysis）方法为服务质量改进的优先级排序提供了分析工具。花青青等（2025）基于SERVQUAL-IPA模型对P银行J支行的金融机构服务质量进行了评价，以重要度为横轴、满意度为纵轴构建四象限矩阵，将指标分为"优先改进区""继续保持区""供给过度区"和"低优先区"\cite{ref15}。该方法的逻辑简明而有力：资源永远是有限的，服务质量改进不应追求"补齐所有短板"，而应聚焦于"高重要度-低满意度"象限中的关键指标。这种以数据驱动优先排序的思路，与DMAIC方法论中的Measure和Analyze阶段高度契合。\par

在服务质量评价的赋权方法方面，AHP（层次分析法）为多准则决策提供了结构化的权重确定框架。Moslem等（2023）在对AHP应用的系统综述中确认，该方法适用于服务质量评价中多维度、多指标的权重分配，但特别强调了两个使用边界：第一，判断矩阵必须依赖于领域专家的真实判断，不得由研究者自行填写；第二，一致性检验（CR < 0.1）是权重有效性的必要条件\cite{ref41}。Nguyen等（2022）在港口服务质量研究中将模糊AHP与IPA组合使用，先用AHP确定各维度的客观权重，再用IPA识别改进优先项，该方法组合为本文第4章提供了跨行业的方法参照\cite{ref42}。需要指出的是，港口服务与小微贷款服务在客户群体、服务流程和评价维度上存在本质差异，该方法组合在金融服务中的适用性尚需结合行业特征进行验证和调整。\par

线上服务质量的维度扩展也为本文的模型改良提供了依据。Raza等（2020）提出的改良e-SERVQUAL模型，在原五维基础上增加了网站设计、隐私安全和客户服务等线上特有维度，其研究证实线上渠道的服务质量维度结构与传统线下渠道存在显著差异\cite{ref43}。田越（2023）在邮储银行J支行的客户满意度研究中，运用AHP对服务质量指标进行了分层赋权，其指标体系的构建逻辑（先分层、后赋权、再排序）对本文的问卷设计具有参考价值\cite{ref44}。\par

综合本节讨论，SERVQUAL提供了测量框架，IPA提供了优先级排序工具，AHP提供了权重确定方法，三者构成了一套完整的服务质量评价组合工具箱。该组合工具在零售银行、线上银行和部分服务业中已有应用，但在小微贷款场景中的系统性应用尚属空白，这正是本文第4章的核心贡献点。\par

\section{2.5 DMAIC与服务流程改进研究}

DMAIC是由精益六西格玛方法论衍生出的结构化改进模型，包含Define（界定）、Measure（测量）、Analyze（分析）、Improve（改进）和Control（控制）五个阶段。该模型的核心逻辑是：以数据驱动的方式识别问题、测量差距、分析根因、设计改进方案并建立持续控制机制，最终实现流程绩效的可测量改善和制度化固化。与一般管理建议或经验驱动的改进方式相比，DMAIC强调"以事实和数据说话"，要求每一个改进决策都建立在可验证的证据基础之上。这一特征使得DMAIC天然适合应用于服务质量改进这样的多触点、多维度的复杂问题场景。\par

精益六西格玛在金融服务领域的应用已经积累了可观的实证证据。Vashishth等（2024）在金融服务领域的实证研究表明，精益六西格玛方法通过减少流程变异和消除浪费，能够显著缩短服务周期并提升客户满意度\cite{ref46}。de Koning、Does和Bisgaard（2008）较早论证了精益六西格玛在金融服务中的适用性，指出金融服务的流程化特征和大量可获取的运营数据为DMAIC方法的落地提供了天然条件\cite{ref47}。Guarraia和Schwedel（2008）从行业实践的角度提出，银行在应用精益六西格玛时面临的主要挑战不在于方法本身，而在于如何将项目目标与客户体验指标有效对齐\cite{ref48}。这一观点对本文具有直接的指导意义：服务质量导向的DMAIC改进，必须将客户感知指标（而不仅是内部效率指标）作为Define阶段的核心输出。\par

Adanigbo等（2021）针对金融科技产品客服中心的运营延迟问题，设计了基于精益六西格玛的改进框架，证实了DMAIC在缩短响应时间、降低投诉升级率方面的显著成效\cite{ref49}。Delahoz-Dominguez等（2024）将六西格玛与DEA（数据包络分析）相结合，构建了银行服务质量评估与改进的双阶段框架：第一阶段用DEA识别效率前沿面，第二阶段用六西格玛方法设计改进路径\cite{ref50}。该研究在评估方法上的探索对本论文第6章的效果评价指标体系设计具有借鉴意义。\par

在国内研究方面，若干学位论文为DMAIC在银行服务改进中的应用提供了场景参考。孙启蒙（2023）以农行T支行为对象，运用精益六西格玛方法对零售业务服务质量进行了系统改善，其DMAIC五阶段的实施路径为本文的第5章框架设计提供了参照\cite{ref12}。唐家驰（2025）基于六西格玛理论对N银行Y分行网点服务管理进行了优化研究，在控制阶段提出了看板监控和月度复盘的制度化方案\cite{ref51}。张磊等（2023）从政策扶持与金融科技协同视角，实证检验了小微企业信贷融资的改善路径，为本研究理解普惠金融服务质量改进的政策与技术背景提供了宏观层面的参照\cite{ref52}。黄玉婷（2023）在M公司的客户服务流程优化研究中，将流程改进与人员培训、绩效激励相结合，这种"流程-人员"协同改进的思路对本论文的方案设计具有启发性\cite{ref53}。\par

值得指出的是，上述国内案例虽然提供了DMAIC在金融服务领域应用的实证，但大多聚焦于银行传统业务或内部流程效率，较少触及小微贷款这一细分场景中的服务质量议题。小微贷款服务的特殊性在于：客户群体的异质性高、服务链条长且多触点、数字化与传统人工服务并存、风险控制与服务效率存在张力。这些特征使得DMAIC在该场景中的应用需要更多的场景适配，而非简单的方法套用。\par

\section{2.6 文献述评与研究空白}

基于以上五条研究脉络的系统梳理，本文归纳出以下三个主要研究缺口。\par

缺口一：小微贷款服务质量闭环研究不足。现有普惠金融和小微贷款文献主要关注融资可得性和融资成本，即"能不能贷到"和"以什么价格贷到"\cite{ref2,ref6}。对于小微贷款客户在申请、审批、放款、贷后、续贷等完整旅程中的服务体验，学术界尚未给予充分关注。服务质量文献虽然提供了成熟的评价理论，但其在小微贷款场景中的应用主要集中在银行零售业务，而非普惠金融机构的小微贷款专项业务\cite{ref14,ref15}。二者的交集领域存在明显的研究空白。\par

缺口二：服务质量评价与DMAIC改进在普惠金融场景中脱节。SERVQUAL及其改良模型解决了"如何评价"的问题，IPA和AHP解决了"如何排序"的问题，DMAIC解决了"如何改进"的问题。然而，在现有研究中，这三种工具往往各自为政：服务质量评价论文通常止步于发现差距和提出建议，缺乏从评价到方案设计的工程化转换；DMAIC应用文献则多从内部运营指标出发，较少将客户感知数据作为改进的驱动变量\cite{ref49,ref12}。将"评价 → 排序 → 改进 → 控制"四个环节串联为一个完整闭环的研究，在普惠金融领域尚付阙如。\par

缺口三：持续控制机制薄弱。即便是为数不多的小微信贷服务改进研究，其重点也集中在问题诊断和方案设计阶段，对于改进方案实施后的制度化保障、指标监控、异常预警和效果持续验证着墨甚少\cite{ref50}。DMAIC的Control阶段强调的不是一次性的效果验证，而是建立防止绩效回弹的常态化机制。在小微贷款服务场景中，这种机制设计的缺失意味着改进效果难以持续，同类服务质量问题容易反复出现。\par

基于上述文献缺口，本文的研究定位如下：以P公司小微贷款业务为案例场景，将改良SERVQUAL（八维32指标）评价、IPA/AHP短板识别和DMAIC改进闭环集成在同一框架中，构建"评价测量 → 优先级排序 → 根因分析 → 工程化改进 → 持续控制"的完整工程管理闭环。这一研究定位既继承并整合了前述五条研究脉络的核心成果，又回应了小微贷款服务质量闭环研究的方法论缺失，为该领域提供了一套可操作的集成性分析框架，为同类普惠金融机构的服务质量改进提供了可操作、可复制的方法框架。\par

\subsection*{理论与方法支撑关系图}

章末以文字形式描述本文的理论与方法支撑关系：

\begin{verbatim}
普惠金融/小微贷款背景 ---> 服务质量问题（第2.1-2.2节）
                               |
                     SERVQUAL测量与期望-感知差距诊断（第2.3-2.4节）
                               |
                     IPA/AHP优先级排序与关键短板识别（第2.4节）
                               |
                     DMAIC根因分析与工程化改进方案（第2.5节）
                               |
                     持续控制与效果评价闭环（第2.5-2.6节）
\end{verbatim}

关系图说明：普惠金融文献（第2.1节）解释了"为什么小微贷款服务质量是值得研究的问题"；金融科技文献（第2.2节）说明了"为什么服务模式的变化使服务质量评价和系统改进变得更加紧迫"；服务质量与满意度文献（第2.3节）建立了"客户感知是可以测量和管理的经营变量"这一前提；SERVQUAL/IPA/AHP文献（第2.4节）提供了测量工具和优先级排序方法；DMAIC文献（第2.5节）提供了从测量到改进再到控制的工程化路径。五条研究脉络在本文中不是平行并列的，而是按"背景 → 变量 → 工具 → 路径 → 目标"的逻辑逐层递进，最终在第2.6节汇合于一个集成的MEM闭环框架。P公司内部数据（第3-6章）为该框架提供了案例验证的土壤。\par

